\chapter{Závěr}
\label{chap:zaver}

Cílem této práce bylo vytvořit komplexní socioekonomický model a interaktivní webovou aplikaci pro analýzu dopadů výstavby velkokapacitních datových center v České republice. V kontextu probíhající revoluce v oblasti umělé inteligence a s ní spojeného extrémního nárůstu poptávky po výpočetním výkonu představuje tato práce unikátní propojení technologického, ekonomického a energetického pohledu na infrastrukturu, která se stává novodobou páteří digitální ekonomiky.

Teoretická část práce poskytla nezbytný základ pro pochopení fungování datových center, jejich klíčových parametrů a globálních trendů, přičemž analýza trhu v USA posloužila jako varovný i inspirativní příklad pro český kontext. Navržený model následně transformoval tyto poznatky do matematických vztahů, které zohledňují specifika českého daňového systému, energetické sítě i environmentálních limitů.

Hlavním zjištěním práce je odhalení velkého paradoxu datových center -- model ukázal, že ačkoliv se jedná o investice v řádech desítek miliard korun, jejich přímý přínos pro národní hospodářství v podobě nových pracovních míst nebo daňových výnosů z mezd je v porovnání s tradičním průmyslem či službami minimální. Práce rovněž identifikovala kritické úzké hrdlo v podobě kapacity přenosové soustavy, která by při nekontrolované výstavbě AI center mohla v energetických špičkách čelit vážným problémům.

Praktickým přínosem práce je funkční webová aplikace a automatizovaný nástroj pro sběr dat (scraper), který zmapoval aktuální stav českého trhu. Aplikace umožňuje nejen odborníkům, ale i tvůrcům veřejných politik a laické veřejnosti simulovat různé scénáře rozvoje a v reálném čase sledovat jejich dopady na spotřebu elektřiny, vody, HDP či státní rozpočet. Tím se podařilo snížit informační bariéru a vytvořit prostor pro věcnou diskusi o budoucnosti české digitální cesty.

Závěrečné vyhodnocení naznačuje, že Česká republika by k podpoře výstavby datových center pro globální technologické giganty měla přistupovat se strategickou opatrností. Samotná fyzická přítomnost infrastruktury bez garantovaného přístupu k výpočetnímu výkonu pro lokální výzkum a vývoj nepřináší státu výraznou přidanou hodnotu. Tato práce tak slouží jako metodický základ pro strategické rozhodování o tom, zda se ČR stane pouze pasivním hostitelem energeticky náročné infrastruktury, nebo aktivním hráčem, který dokáže výpočetní výkon využít k posílení vlastní technologické suverenity a konkurenceschopnosti.


